\section{September 8, 2026}

We continue the construction of measures from premeasures. Recall that
if $\mu_0$ is a premeasure on an algebra $\mathcal A$, then
\[
  \mu^*(E)
  =\inf\left\{
    \sum_{j=1}^{\infty}\mu_0(A_j):
    A_j\in\mathcal A,\quad
    E\subseteq\bigcup_{j=1}^{\infty}A_j
  \right\}
\]
defines an outer measure. Every member of $\mathcal A$ is
$\mu^*$-measurable, and
\[
  \mu=\left.\mu^*\right|_{\mathcal M(\mathcal A)}
\]
is a measure extending $\mu_0$.

\subsection{Uniqueness of the extension}

The extension is unique on sets of finite measure. Under a
$\sigma$-finiteness assumption, it is unique everywhere.

\begin{theorem}[Uniqueness of the premeasure extension]
  \label{thm:uniqueness-premeasure-extension}
  Let $\mu_0$ be a premeasure on an algebra $\mathcal A$, and let
  $\mu$ be the measure on $\mathcal M(\mathcal A)$ obtained from the
  induced outer measure. Suppose that $\nu$ is another measure on
  $\mathcal M(\mathcal A)$ extending $\mu_0$.
  \begin{enumerate}[label=(\alph*)]
    \item If $E\in\mathcal M(\mathcal A)$ and $\mu(E)<\infty$, then
      $\nu(E)=\mu(E)$.
    \item If $\mu_0$ is $\sigma$-finite, then $\nu=\mu$ on
      $\mathcal M(\mathcal A)$.
  \end{enumerate}
\end{theorem}

\begin{proof}
  We divide the proof of part (a) into three steps.

  First, let $E\in\mathcal M(\mathcal A)$ and suppose that
  $E\subseteq\bigcup_{j=1}^{\infty}A_j$, where $A_j\in\mathcal A$.
  Since $\nu$ extends $\mu_0$, countable subadditivity gives
  \[
    \nu(E)
    \leq\nu\left(\bigcup_{j=1}^{\infty}A_j\right)
    \leq\sum_{j=1}^{\infty}\nu(A_j)
    =\sum_{j=1}^{\infty}\mu_0(A_j).
  \]
  Taking the infimum over all such covers yields
  \begin{equation}\label{eq:nu-below-mu-sep-08}
    \nu(E)\leq\mu^*(E)=\mu(E).
  \end{equation}

  Second, suppose that
  \[
    A=\bigcup_{j=1}^{\infty}A_j,
    \qquad A_j\in\mathcal A.
  \]
  Set $B_n=\bigcup_{j=1}^nA_j$. Then $B_n\in\mathcal A$ and
  $B_n\uparrow A$. Continuity from below, together with the fact that
  $\mu$ and $\nu$ both extend $\mu_0$, gives
  \[
    \nu(A)
    =\lim_{n\to\infty}\nu(B_n)
    =\lim_{n\to\infty}\mu(B_n)
    =\mu(A).
  \]

  Third, let $E\in\mathcal M(\mathcal A)$ satisfy $\mu(E)<\infty$,
  and fix $\varepsilon>0$. By the definition of $\mu^*$, there are
  $A_j\in\mathcal A$ such that, with
  $A=\bigcup_{j=1}^{\infty}A_j$,
  \[
    E\subseteq A,
    \qquad
    \mu(A)\leq\sum_{j=1}^{\infty}\mu_0(A_j)
      <\mu(E)+\varepsilon.
  \]
  In particular, $\mu(A\setminus E)<\varepsilon$. By the second step
  and \eqref{eq:nu-below-mu-sep-08},
  \begin{align*}
    \mu(E)
    &\leq\mu(A)
     =\nu(A)\\
    &=\nu(E)+\nu(A\setminus E)\\
    &\leq\nu(E)+\mu(A\setminus E)
     <\nu(E)+\varepsilon.
  \end{align*}
  Letting $\varepsilon\downarrow0$ shows that
  $\mu(E)\leq\nu(E)$. Together with
  \eqref{eq:nu-below-mu-sep-08}, this proves part (a).

  For part (b), $\sigma$-finiteness provides sets
  $A_j\in\mathcal A$ such that
  \[
    X=\bigcup_{j=1}^{\infty}A_j,
    \qquad
    \mu_0(A_j)<\infty.
  \]
  By replacing $A_j$ with
  $A_j\setminus\bigcup_{i=1}^{j-1}A_i$, we may assume that the sets are
  pairwise disjoint. If $E\in\mathcal M(\mathcal A)$, then
  $\mu(E\cap A_j)\leq\mu_0(A_j)<\infty$, so part (a) applies to each
  $E\cap A_j$. Therefore
  \[
    \mu(E)
    =\sum_{j=1}^{\infty}\mu(E\cap A_j)
    =\sum_{j=1}^{\infty}\nu(E\cap A_j)
    =\nu(E).
  \]
\end{proof}

\subsection{Lebesgue--Stieltjes measures on \texorpdfstring{$\bb{R}$}{R}}

Let $\mu$ be a finite Borel measure on $\bb{R}$. Its
\emph{distribution function} is
\begin{equation}\label{eq:distribution-function-sep-08}
  F(x)=\mu((-\infty,x]),
  \qquad x\in\bb{R}.
\end{equation}
This function records the measure of every left half-line.

\begin{proposition}\label{prop:distribution-function-properties-sep-08}
  The distribution function $F$ has the following properties.
  \begin{enumerate}[label=(\roman*)]
    \item It is increasing.
    \item It is right-continuous.
    \item For every $a<b$,
      \[
        \mu((a,b])=F(b)-F(a).
      \]
  \end{enumerate}
\end{proposition}

\begin{proof}
  If $x\leq y$, then $(-\infty,x]\subseteq(-\infty,y]$, so
  monotonicity of $\mu$ gives $F(x)\leq F(y)$.

  If $x_n\downarrow x$, then
  \[
    (-\infty,x_n]\downarrow(-\infty,x].
  \]
  Since $\mu$ is finite, continuity from above gives
  \[
    \lim_{n\to\infty}F(x_n)
    =\lim_{n\to\infty}\mu((-\infty,x_n])
    =\mu((-\infty,x])
    =F(x).
  \]
  Hence $F$ is right-continuous. Finally,
  \[
    (-\infty,b]=(-\infty,a]\mathbin{\dot\cup}(a,b],
  \]
  and finite additivity gives
  \[
    \mu((a,b])
    =\mu((-\infty,b])-\mu((-\infty,a])
    =F(b)-F(a).
  \]
\end{proof}

\begin{example}[Distribution function of a Dirac measure]
  Let $\mu=\delta_{x_0}$ be the Dirac measure at $x_0\in\bb{R}$.
  Then
  \[
    F(x)=\delta_{x_0}((-\infty,x])
    =\begin{cases}
      0,&x<x_0,\\
      1,&x\geq x_0.
    \end{cases}
  \]
  Thus an atom of a measure appears as a jump in its distribution
  function.
\end{example}

\begin{remark}[Measures finite on bounded sets]
  A Borel measure need not be finite on all of $\bb{R}$. If $\mu$ is
  finite on every bounded Borel set, define the normalized distribution
  function
  \[
    \widetilde F(x)
    =\begin{cases}
      \mu((0,x]),&x>0,\\
      0,&x=0,\\
      -\mu((x,0]),&x<0.
    \end{cases}
  \]
  Then $\widetilde F$ is increasing and right-continuous, and
  \[
    \mu((a,b])=\widetilde F(b)-\widetilde F(a)
  \]
  for all $a<b$. Adding a constant to $\widetilde F$ does not change
  these interval increments.
\end{remark}

We now reverse the preceding construction at the level of the algebra of
half-open intervals. Let $F\colon\bb{R}\to\bb{R}$ be increasing and
right-continuous, and set
\[
  F(-\infty)=\lim_{x\to-\infty}F(x),
  \qquad
  F(\infty)=\lim_{x\to\infty}F(x),
\]
where these limits are allowed to be infinite. Let $\mathcal A$ be the
algebra of finite unions of intervals of the forms
\[
  (a,b],\qquad (-\infty,b],\qquad (a,\infty),
\]
together with $\varnothing$. Equivalently, its members are finite unions
of intervals $(a,b]\cap\bb{R}$ with
$-\infty\leq a<b\leq\infty$.

\begin{proposition}[The Lebesgue--Stieltjes premeasure]
  \label{prop:lebesgue-stieltjes-premeasure-sep-08}
  Suppose that
  \[
    E=\bigcup_{j=1}^N(a_j,b_j]
  \]
  is a representation of $E\in\mathcal A$ as a finite union of pairwise
  disjoint half-open intervals, with extended endpoints allowed as
  above. Define
  \begin{equation}\label{eq:stieltjes-premeasure-sep-08}
    \mu_0(E)
    =\sum_{j=1}^N\bigl(F(b_j)-F(a_j)\bigr),
    \qquad
    \mu_0(\varnothing)=0.
  \end{equation}
  Then $\mu_0\colon\mathcal A\to[0,\infty]$ is well defined and is a
  premeasure.
\end{proposition}

\begin{proof}
  We first show that the value in
  \eqref{eq:stieltjes-premeasure-sep-08} is independent of the chosen
  representation. Suppose that a half-open interval is partitioned into
  finitely many pairwise disjoint half-open intervals:
  \[
    (a,b]=\bigcup_{j=1}^N(a_j,b_j].
  \]
  After relabeling, their endpoints satisfy
  \[
    a=a_1<b_1=a_2<b_2=\cdots<b_{N-1}=a_N<b_N=b.
  \]
  Hence the sum telescopes:
  \[
    \sum_{j=1}^N\bigl(F(b_j)-F(a_j)\bigr)=F(b)-F(a).
  \]
  The same argument applies when an endpoint is infinite, using the
  conventions for $F(-\infty)$ and $F(\infty)$.

  Now suppose that
  \[
    \bigcup_{i=1}^N I_i=\bigcup_{j=1}^M J_j,
  \]
  where both families consist of pairwise disjoint half-open intervals.
  The nonempty intersections $I_i\cap J_j$ partition both sides. By the
  preceding telescoping identity,
  \begin{align*}
    \sum_{i=1}^N\mu_0(I_i)
    &=\sum_{i=1}^N\sum_{j=1}^M\mu_0(I_i\cap J_j)\\
    &=\sum_{j=1}^M\sum_{i=1}^N\mu_0(I_i\cap J_j)
     =\sum_{j=1}^M\mu_0(J_j).
  \end{align*}
  Thus $\mu_0$ is well defined; the same calculation also gives finite
  additivity.

  It remains to prove countable additivity. After decomposing all sets
  involved into their interval components and intersecting with each
  component of the union, it is enough to consider pairwise disjoint
  half-open intervals $I_j$ such that
  \[
    (a,b]=\bigcup_{j=1}^{\infty}I_j.
  \]
  Finite additivity and monotonicity give, for every $n$,
  \[
    \mu_0((a,b])
    \geq\mu_0\left(\bigcup_{j=1}^n I_j\right)
    =\sum_{j=1}^n\mu_0(I_j).
  \]
  Letting $n\to\infty$ yields
  \begin{equation}\label{eq:stieltjes-first-inequality-sep-08}
    \mu_0((a,b])\geq\sum_{j=1}^{\infty}\mu_0(I_j).
  \end{equation}

  For the reverse inequality, first suppose that $a,b\in\bb{R}$ and
  write $I_j=(a_j,b_j]$. Fix $\varepsilon>0$. Right-continuity provides
  $0<\delta<b-a$ and $\delta_j>0$ such that
  \[
    F(a+\delta)-F(a)<\varepsilon,
    \qquad
    F(b_j+\delta_j)-F(b_j)<\varepsilon 2^{-j}.
  \]
  The open intervals $(a_j,b_j+\delta_j)$ cover the compact interval
  $[a+\delta,b]$, so finitely many of them suffice. After discarding
  redundant intervals, relabel this finite subcover from left to right
  with indices $1,\ldots,r$. Its total enlargement still satisfies
  \[
    \sum_{j=1}^r\bigl(F(b_j+\delta_j)-F(b_j)\bigr)<\varepsilon,
  \]
  and the ordering may be chosen so that
  \[
    a_1<a+\delta,\qquad
    a_{j+1}<b_j+\delta_j\ (1\leq j<r),\qquad
    b<b_r+\delta_r.
  \]
  Monotonicity of $F$ then gives
  \begin{align*}
    F(b)-F(a)
    &<F(b)-F(a+\delta)+\varepsilon\\
    &\leq F(b_r+\delta_r)-F(a_1)+\varepsilon\\
    &=F(b_r+\delta_r)-F(a_r)
      +\sum_{j=1}^{r-1}\bigl(F(a_{j+1})-F(a_j)\bigr)
      +\varepsilon\\
    &\leq\sum_{j=1}^r
      \bigl(F(b_j+\delta_j)-F(a_j)\bigr)+\varepsilon\\
    &<\sum_{j=1}^r\bigl(F(b_j)-F(a_j)\bigr)+2\varepsilon\\
    &\leq\sum_{j=1}^{\infty}\mu_0(I_j)+2\varepsilon.
  \end{align*}
  Letting $\varepsilon\downarrow0$ proves the reverse of
  \eqref{eq:stieltjes-first-inequality-sep-08} for bounded intervals.

  If the interval on the left is unbounded, apply the bounded case to
  its truncations. For example,
  \[
    (-m,b]\uparrow(-\infty,b],
    \qquad
    \mu_0((-m,b])=F(b)-F(-m)\uparrow F(b)-F(-\infty),
  \]
  while the bounded case gives
  \[
    \mu_0((-m,b])
    =\sum_{j=1}^{\infty}\mu_0\bigl(I_j\cap(-m,b]\bigr)
    \leq\sum_{j=1}^{\infty}\mu_0(I_j).
  \]
  Letting $m\to\infty$ proves the desired inequality for
  $(-\infty,b]$. The same argument uses
  $(a,m]\uparrow(a,\infty)$ or $(-m,m]\uparrow\bb{R}$ for the other
  unbounded cases. Therefore
  \[
    \mu_0\left(\bigcup_{j=1}^{\infty}I_j\right)
    =\sum_{j=1}^{\infty}\mu_0(I_j),
  \]
  whenever the union belongs to $\mathcal A$, so $\mu_0$ is a
  premeasure.
\end{proof}

The premeasure extension theorem from the preceding lecture therefore
applies to $\mu_0$ on $\mathcal A$.

\begin{remark}[Two Lebesgue--Stieltjes examples]
  The Borel measure determined by the increments
  $F(b)-F(a)$ is denoted by $\mu_F$, and its completion is called the
  \emph{Lebesgue--Stieltjes measure} associated with $F$.
  \begin{enumerate}[label=(\roman*)]
    \item For $F(x)=x$,
      \[
        \mu_F((a,b])=b-a.
      \]
      The completed measure is Lebesgue measure on $\bb{R}$.
    \item For the right-continuous step function
      \[
        F(x)=\mathbf{1}_{[x_0,\infty)}(x),
      \]
      the associated measure is the Dirac measure
      $\mu_F=\delta_{x_0}$.
  \end{enumerate}
\end{remark}
